\section{はじめに}
\subsection{背景}
バスサイネージシステムは，バスの到着予定時刻や運行状況をリアルタイムに表示する電子掲示板で，利用者にとって有用な情報源となっている．
しかし，多くのシステムは個々の事情に合わせてオーダーメイドで作成され，完成度は高いものの，導入には大きなコストがかかる．
そのため，公民館や個人商店など，導入コストの負担が難しい施設においては，バスサイネージシステムが設置されてこなかった．
そこで，小規模なバス停においてもバスサイネージシステムを導入できるようにするため，「設置者自身による設定」によって動作する新たなバスサイネージシステムを構築する．

バスの運行情報を標準化された形式で取得できるGTFSが普及しつつある~\cite{GTFS}．
日本国内においては，国土交通省が拡張・策定しているGTFS-JPが広く使われている~\cite{GTFSJPv4}．
2026年7月現在，GTFS形式による公共交通オープンデータは全国800事業者において整備されている~\cite{GTFSJPOpenDataMap}．
また，リアルタイム情報を提供するGTFS-RTについても整備が進んでいる~\cite{GTFSJPv4}．
このGTFSを入力とすることで，GTFSを基盤とした汎用バスサイネージシステムを作成できる可能性がある．

\subsection{これまでの研究}
前回研究では，GTFS-JPを用いた設置者自身で設定可能なバスサイネージシステムを構築した．
同研究では，GTFSから取り出せる基本的な運行情報を利用した表示を実現した．
一方で，設置者が表示対象等を設定できることと，その限界が明らかになった．
また，GTFSから利用者向けの表示に必要な意味情報を取得する部分が課題となった~\cite{oai:ipsj.ixsq.nii.ac.jp:02007105}．

\subsection{本研究の目的と位置づけ}
本研究では，GTFSから取得したバス運行データをバス利用者向けのバスサイネージ表示に適した情報へ自動的に変換するため，Claude Code\footnote{\url{https://code.claude.com/docs/}}やCodex\footnote{\url{https://openai.com/codex/}}のようなAIエージェントを用いた設定支援手法を提案する．
GTFSは運行データとしては構造化されている一方，バス利用者向け向け案内情報の意味構造は統一されていない．
GTFSをそのまま利用するだけでは得られない情報をAIエージェントにより抽出することで，セルフサービスでの設置に必要な設定作業を削減する．
